% Options for packages loaded elsewhere
\PassOptionsToPackage{unicode}{hyperref}
\PassOptionsToPackage{hyphens}{url}
\documentclass[
  ignorenonframetext,
]{beamer}
\newif\ifbibliography
\usepackage{pgfpages}
\setbeamertemplate{caption}[numbered]
\setbeamertemplate{caption label separator}{: }
\setbeamercolor{caption name}{fg=normal text.fg}
\beamertemplatenavigationsymbolsempty
% remove section numbering
\setbeamertemplate{part page}{
  \centering
  \begin{beamercolorbox}[sep=16pt,center]{part title}
    \usebeamerfont{part title}\insertpart\par
  \end{beamercolorbox}
}
\setbeamertemplate{section page}{
  \centering
  \begin{beamercolorbox}[sep=12pt,center]{section title}
    \usebeamerfont{section title}\insertsection\par
  \end{beamercolorbox}
}
\setbeamertemplate{subsection page}{
  \centering
  \begin{beamercolorbox}[sep=8pt,center]{subsection title}
    \usebeamerfont{subsection title}\insertsubsection\par
  \end{beamercolorbox}
}
% Prevent slide breaks in the middle of a paragraph
\widowpenalties 1 10000
\raggedbottom
\AtBeginPart{
  \frame{\partpage}
}
\AtBeginSection{
  \ifbibliography
  \else
    \frame{\sectionpage}
  \fi
}
\AtBeginSubsection{
  \frame{\subsectionpage}
}
\usepackage{iftex}
\ifPDFTeX
  \usepackage[T1]{fontenc}
  \usepackage[utf8]{inputenc}
  \usepackage{textcomp} % provide euro and other symbols
\else % if luatex or xetex
  \usepackage{unicode-math} % this also loads fontspec
  \defaultfontfeatures{Scale=MatchLowercase}
  \defaultfontfeatures[\rmfamily]{Ligatures=TeX,Scale=1}
\fi
\usepackage{lmodern}
\ifPDFTeX\else
  % xetex/luatex font selection
\fi
% Use upquote if available, for straight quotes in verbatim environments
\IfFileExists{upquote.sty}{\usepackage{upquote}}{}
\IfFileExists{microtype.sty}{% use microtype if available
  \usepackage[]{microtype}
  \UseMicrotypeSet[protrusion]{basicmath} % disable protrusion for tt fonts
}{}
\makeatletter
\@ifundefined{KOMAClassName}{% if non-KOMA class
  \IfFileExists{parskip.sty}{%
    \usepackage{parskip}
  }{% else
    \setlength{\parindent}{0pt}
    \setlength{\parskip}{6pt plus 2pt minus 1pt}}
}{% if KOMA class
  \KOMAoptions{parskip=half}}
\makeatother
\usepackage{longtable,booktabs,array}
\newcounter{none} % for unnumbered tables
\usepackage{calc} % for calculating minipage widths
\usepackage{caption}
% Make caption package work with longtable
\makeatletter
\def\fnum@table{\tablename~\thetable}
\makeatother
\usepackage{graphicx}
\makeatletter
\newsavebox\pandoc@box
\newcommand*\pandocbounded[1]{% scales image to fit in text height/width
  \sbox\pandoc@box{#1}%
  \Gscale@div\@tempa{\textheight}{\dimexpr\ht\pandoc@box+\dp\pandoc@box\relax}%
  \Gscale@div\@tempb{\linewidth}{\wd\pandoc@box}%
  \ifdim\@tempb\p@<\@tempa\p@\let\@tempa\@tempb\fi% select the smaller of both
  \ifdim\@tempa\p@<\p@\scalebox{\@tempa}{\usebox\pandoc@box}%
  \else\usebox{\pandoc@box}%
  \fi%
}
% Set default figure placement to htbp
\def\fps@figure{htbp}
\makeatother
\setlength{\emergencystretch}{3em} % prevent overfull lines
\providecommand{\tightlist}{%
  \setlength{\itemsep}{0pt}\setlength{\parskip}{0pt}}
% Slide layout defaults for warning-clean scientific decks.
\usepackage{etoolbox}
\IfFileExists{xurl.sty}{\usepackage{xurl}}{}
\IfFileExists{seqsplit.sty}{\usepackage{seqsplit}}{\newcommand{\seqsplit}[1]{#1}}
\protected\def\breakseq#1{\seqsplit{#1}}
\protected\def\breaktt#1{\begingroup\ttfamily\seqsplit{#1}\endgroup}
\setlength{\emergencystretch}{6em}
\tolerance=5000
\hbadness=10000
\hfuzz=1pt
\setlength{\tabcolsep}{2pt}
\AtBeginEnvironment{longtable}{\tiny\renewcommand{\arraystretch}{0.86}\setlength{\tabcolsep}{1pt}}
\AtBeginEnvironment{tabular}{\tiny\renewcommand{\arraystretch}{0.86}\setlength{\tabcolsep}{1pt}}

% Natbib and cross-reference fallbacks — slides don't load natbib
% or cleveref, but combined-PDF manuscript prose may emit these
% commands. The fallback renders citations as a bracketed key list
% and cross-references as detokenized labels so slides don't fail on
% undefined control sequences or raw underscores. \providecommand is
% a no-op if packages load later.
\providecommand{\citep}[1]{[#1]}
\providecommand{\citet}[1]{#1}
\providecommand{\citealp}[1]{#1}
\providecommand{\citeauthor}[1]{#1}
\providecommand{\citeyear}[1]{#1}
\providecommand{\cref}[1]{\texttt{\detokenize{#1}}}
\providecommand{\Cref}[1]{\texttt{\detokenize{#1}}}

\newtheorem{proposition}{Proposition}
\newtheorem{hypothesis}{Hypothesis}
\usepackage{bookmark}
\IfFileExists{xurl.sty}{\usepackage{xurl}}{} % add URL line breaks if available
\urlstyle{same}
% fallback for those not using the hyperref driver hyperxmp:
\makeatletter
\@ifundefined{xmpquote}{\newcommand{\xmpquote}[1]{#1}}{}
\makeatother
\hypersetup{
  hidelinks,
  pdfcreator={LaTeX via pandoc}}

\author{\texorpdfstring{}{}}
\date{}

\begin{document}

\section{Deviation register}\label{sec:deviations}

\begin{frame}[fragile,allowframebreaks]{Deviation register}
A registered report must state, explicitly, every way the executed
analysis departed from the frozen plan. This template treats that
register as a first-class artifact: \breaktt{build\_deviation\_ledger}
produces one row per executed outcome and model, classifying each as
\texttt{ok} (registered), \texttt{warning} (an unregistered element
carried with a documented rationale), or \texttt{error} (an unregistered
element with no rationale). \breaktt{compare\_analysis\_to\_registration}
downgrades a documented model change from \texttt{error} to
\texttt{warning} and refuses to accept any deviation whose
\texttt{rationale} is empty.
\end{frame}

\begin{frame}[fragile,allowframebreaks]{Demonstration deviations}
\protect\phantomsection\label{demonstration-deviations}
To exercise the machinery, the demonstration executes a deliberately
plan-divergent analysis: it adds a \texttt{secondary\_score} endpoint
and swaps the registered \breaktt{permutation\_test} for a
\texttt{linear\_model}, each accompanied by a rationale. The resulting
ledger is:

{\def\LTcaptype{none} % do not increment counter
\begin{longtable}[]{@{}llllll@{}}
\toprule\noalign{}
Order & Kind & Target & Registered & Documented & Severity \\
\midrule\noalign{}
\endhead
0 & outcome & \texttt{primary\_score} & yes & --- & \texttt{ok} \\
1 & outcome & \texttt{secondary\_score} & no & yes & \texttt{warning} \\
2 & model & \texttt{linear\_model} & no & yes & \texttt{warning} \\
\bottomrule\noalign{}
\end{longtable}
}

Because both departures are documented, the review packet remains valid:
it reports \texttt{primary\_score} as the only confirmatory outcome,
\texttt{secondary\_score} as exploratory, and an integrity score of
\texttt{0.9} (a single \texttt{0.10} warning penalty against a perfect
plan). Had either deviation lacked a rationale, the same code would have
raised a blocking \breaktt{deviation\_without\_rationale} or
\breaktt{unregistered\_outcome} error, and the packet would be invalid.
\breakseq{[@fig:deviation\_timeline]} plots this register along the
registered-to-executed axis.

\begin{figure}
\centering
\pandocbounded{\includegraphics[keepaspectratio,alt={Preregistration-versus-executed deviation ledger timeline, rendered by plot\_deviation\_timeline from build\_deviation\_ledger and deviation\_timeline. Each marker is a ledger row coloured by severity: the registered primary\_score outcome is green (ok), while the added secondary\_score endpoint and the linear\_model substitution are amber (warning, documented deviations). No marker is red because every departure carries a rationale. An undocumented departure would appear in red and invalidate the review packet.}]{../figures/deviation_timeline.png}}
\caption{Preregistration-versus-executed deviation ledger timeline,
rendered by \protect\breaktt{plot\_deviation\_timeline} from
\protect\breaktt{build\_deviation\_ledger} and
\protect\breaktt{deviation\_timeline}. Each marker is a ledger row
coloured by severity: the registered \protect\texttt{primary\_score}
outcome is green (\protect\texttt{ok}), while the added
\protect\texttt{secondary\_score} endpoint and the
\protect\texttt{linear\_model} substitution are amber
(\protect\texttt{warning}, documented deviations). No marker is red
because every departure carries a rationale. An undocumented departure
would appear in red and invalidate the review
packet.}\label{fig:deviation_timeline}
\end{figure}
\end{frame}

\begin{frame}[fragile,allowframebreaks]{Reading the register}
\protect\phantomsection\label{reading-the-register}
The register is the mechanism that keeps {[}Section @sec:results{]}
honest: the confirmatory claim is restricted to the registered
\texttt{primary\_score} outcome, and the exploratory
\texttt{secondary\_score} endpoint is visible, labelled, and excluded
from confirmatory interpretation. Forks should extend this ledger ---
never edit the frozen registration --- whenever the executed analysis
must diverge from plan.
\end{frame}

\end{document}
